
%% $RCSfile: proj_proposal.tex,v $
%% $Revision: 1.4 $
%% $Date: 2017/10/06 02:55:50 $
%% $Author: kevin $

\documentclass[11pt, a4paper, twoside, openright]{article}

\usepackage{parskip}

\usepackage{float} % lets you have non-floating floats

\usepackage{url} % for typesetting urls

\usepackage[acronym]{glossaries} % for better handling of acronyms

\usepackage{pgfgantt} % gantt chart

\raggedbottom

%% Example acronyms used in this document. 
\newacronym{hci}{HCI}{Human-Computer Interaction}
\newacronym{acm}{ACM}{Association for Computing Machinery}
\newacronym{ieee}{IEEE}{Institute of Electrical and Electronics Engineers}

%% We don't want figures to float so we define
%%
\newfloat{fig}{thp}{lof}[section]
\floatname{fig}{Figure}

%% These are standard LaTeX definitions for the document
%%
\title{Scotland Yard Multiplayer Game as Map Overlay}
\author{Ming Bao}

%% This file can be used for creating a wide range of reports
%%  across various Schools
%%
%% Set up some things, mostly for the front page, for your specific document
%
% Current options are:
% [ecs|msor|sms]          Which school you are in.
%                         (msor option retained for reproducing old data)
% [bschonscomp|mcompsci]  Which degree you are doing
%                          You can also specify any other degree by name
%                          (see below)
% [font|image]            Use a font or an image for the VUW logo
%                          The font option will only work on ECS systems
%
\usepackage[image,ecs]{vuwproject} 

% You should specifiy your supervisor here with
%     \supervisor{Firstname Lastname}
% use \supervisors if there are more than one supervisor
\supervisors{Jens Dietrich, Stuart Marshall}

% Unless you've used the bschonscomp or mcompsci
%  options above use
%   \otherdegree{OTHER DEGREE OR DIPLOMA NAME}
% here to specify degree
\otherdegree{ENGR489}

% Comment this out if you want the date printed.
\date{}

\begin{document}
	
	% Make the page numbering roman, until after the contents, etc.
	\frontmatter
	
	%%%%%%%%%%%%%%%%%%%%%%%%%%%%%%%%%%%%%%%%%%%%%%%%%%%%%%%
	
	\begin{abstract}
	This project develops a web-based multiplayer version of Scotland Yard overlaid on a real-world map of the Wellington region. Physical gameplay limits accessibility, and existing digital versions lack local geographic integration. The system will implement core game logic and user interface components within a browser environment, integrating real-time multiplayer communication and interactive mapping technologies to support location-based gameplay. The expected outcome is a functional online platform that enables synchronous play, enforces game rules automatically, and visualises movement across the Wellington map. The project will be evaluated through functional testing of game mechanics, performance testing of multiplayer interactions, and user testing to assess usability and overall gameplay experience. The completed system will also be suitable for demonstration at university Open Day events and for promotional use, showcasing the level of programming skills that can be achieved through interactive gameplay.
	\end{abstract}
	
	%%%%%%%%%%%%%%%%%%%%%%%%%%%%%%%%%%%%%%%%%%%%%%%%%%%%%%%
	\maketitle
	%\tableofcontents
	
	% we want a list of the figures we defined
	%\listof{fig}{Figures}
	
	%%%%%%%%%%%%%%%%%%%%%%%%%%%%%%%%%%%%%%%%%%%%%%%%%%%%%%%
	
	\mainmatter
	
	%\setlength{\parindent}{0pt}
	%\setlength{\parskip}{4pt}
	
	%%%%%%%%%%%%%%%%%%%%%%%%%%%%%%%%%%%%%%%%%%%%%%%%%%%%%%%
	
	\section{Introduction}
	
	Many board games are limited by physical components and the need for players to be in the same location. Scotland Yard, a hidden-movement strategy game based on movement across a map, is well suited to a digital implementation but most existing versions simply replicate the original board and do not use modern web or mapping technologies.
	
	This project aims to implement a web-based multiplayer version of Scotland Yard that overlays game logic and user interface elements onto a real-world map of the Wellington region using a map API. A browser-based frontend will provide an interactive interface for players, while a backend service will manage game state, enforce rules, and synchronise player actions in real time. The system is also intended to be suitable for demonstration in university Open Day events and promotional settings, showcasing interactive web technologies and real-time multiplayer system design.
	
	The system will be evaluated through functional testing of the game mechanics, performance testing of multiplayer interactions, and user testing to assess usability and gameplay experience. The project will primarily require software resources including web development frameworks, mapping APIs, and networking tools for real-time communication.
	
	\section{The Problem}
	
	Board games such as Scotland Yard are traditionally limited to physical play, requiring co-located participants and a fixed board layout. Although digital versions exist, they typically replicate the original board and do not integrate real-world geographic data or modern web-based multiplayer capabilities. There is an opportunity to extend the game into a location-aware, browser-based platform that combines structured game logic with interactive mapping and contemporary user interface design.
	
	The engineering problem is to implement a web-based multiplayer system that enforces the rules of Scotland Yard while overlaying gameplay onto a real-world map of the Wellington region using an API from a mapping service. This requires integrating turn-based game logic, real-time client--server communication, and server push mechanisms (such as WebSockets) to ensure timely updates of game state across all clients, along with geographic data processing and a responsive, intuitive frontend. The system must synchronise player actions reliably, ensure fairness through programmatic rule enforcement, and present clear visual feedback through effective interface design. In addition, the system should be robust and intuitive enough to be demonstrated to a general audience during Open Day events and used as part of university promotional activities.
	
	The aim of this project is to implement and evaluate an online multiplayer version of Scotland Yard that runs entirely in a web browser, uses a map API to render Wellington as the game board, provides a well-designed frontend for interaction and visualisation, and supports real-time multiplayer gameplay. The project will demonstrate correctness of game mechanics, robustness of multiplayer synchronisation, usability through structured testing and user evaluation, and suitability for public demonstration.
	
	\section{Proposed Solution}
	
	The project will be implemented as a web-based multiplayer system that integrates game logic, interactive mapping, and a responsive user interface. The system will follow a client–server architecture where the backend manages the game state and enforces the rules of Scotland Yard, while the frontend provides an interactive interface for players within a web browser.
	
	The backend will implement the core game engine, including player roles, turn management, movement rules, ticket usage, and win conditions. It will also manage game sessions and synchronise actions between players using real-time communication technologies. This ensures that all players share a consistent game state and that the rules of the game are enforced automatically.
	
	The frontend will be responsible for presenting the game visually and allowing players to interact with it. An interactive map API will be used to display the Wellington region as the game board, with game locations mapped to nodes on the map. The interface will allow players to view available moves, select transport options, and track the progress of the game. Particular attention will be given to designing a clear and responsive user interface so that players can easily understand the game state and perform actions.
	
	The development process will proceed iteratively. Initial stages will focus on designing the system architecture, selecting appropriate technologies, and producing UI prototypes. Implementation will then focus on the core game logic, followed by integration with the map interface and multiplayer communication features. Subsequent stages will refine the system through testing, usability improvements, and deployment preparation.
	
	The project will be carried out over a 24-week period and will follow an iterative development process with clearly defined milestones. Early stages will focus on planning and design to reduce technical risk before implementation begins.
	
	Project planning and design (Weeks 5–9).
	Define system requirements and architecture, select a suitable map API, and design the overall system structure. Produce frontend interface designs and backend architecture diagrams describing the game logic, multiplayer communication, and data flow.
	Output: system architecture diagrams, UI mockups, and technology selection documentation.
	
	Core implementation (Weeks 9–15).
	Implement the backend game engine that enforces Scotland Yard rules, manages game sessions, and handles multiplayer communication. Develop the browser-based frontend and integrate the mapping API to render the Wellington map and player movement.
	Output: a working prototype supporting multiplayer gameplay and map-based interaction.
	
	Testing and evaluation (Weeks 15–18).
	Perform functional testing of game rules, test multiplayer synchronisation and performance, and conduct user testing to evaluate usability and gameplay experience.
	Output: test results and evaluation report.
	
	Refinement and deployment (Weeks 18–24).
	Improve the interface, optimise performance, and deploy the system with appropriate hosting and infrastructure.
	Output: final deployed system and completed project report.
	
	
	\resizebox{\textwidth}{!}{
		\begin{ganttchart}[
			y unit title=0.7cm,
			y unit chart=0.6cm,
			vgrid,
			hgrid
			]{1}{24}
			
			\gantttitle{Project Timeline (Weeks)}{24} \\
			\gantttitlelist{1,...,24}{1} \\
			
			\ganttgroup{Planning}{5}{9} \\
			\ganttbar{Project planning}{5}{6} \\
			\ganttbar{System architecture design}{6}{8} \\
			\ganttbar{UI design and prototyping}{6}{9} \\
			\ganttbar{Map API selection}{7}{9} \\
			
			\ganttgroup{Development}{9}{15} \\
			\ganttbar{Backend game logic implementation}{9}{13} \\
			\ganttbar{Frontend development}{10}{14} \\
			\ganttbar{Map API integration}{11}{14} \\
			\ganttbar{Working multiplayer prototype}{14}{15} \\
			
			\ganttgroup{Testing and Evaluation}{15}{18} \\
			\ganttbar{Functional testing}{15}{17} \\
			\ganttbar{User testing and evaluation}{17}{18} \\
			
			\ganttgroup{Refinement and Deployment}{18}{24} \\
			\ganttbar{System improvements and optimisation}{18}{22} \\
			\ganttbar{Hosting and infrastructure setup}{20}{23} \\
			\ganttbar{Final report and documentation}{21}{24}
			
		\end{ganttchart}
	}
	
	There is two two week mid trimester breaks and one five week trimester break that I could use to catch up if I'm behind at any point.
	
	\section{Evaluating your Solution}
	
	The completed system will be evaluated using a combination of automated testing, performance testing, and user evaluation. These approaches will assess the correctness of the implementation, the reliability of the multiplayer system, and the usability of the interface.
	
	A combination of \textbf{white-box} and \textbf{black-box testing} techniques will be applied. White-box testing will focus on internal logic and code paths, while black-box testing will validate system behaviour from the user’s perspective without knowledge of the implementation.
	
	\textbf{Unit testing} will be used to verify the correctness of the backend game logic. Individual components such as player movement, ticket usage, turn management, and win conditions will be tested in isolation to ensure they behave as expected. Automated tests will confirm that invalid moves are rejected and that the game state remains consistent after each action.
	
	\textbf{Mock testing} will be used to simulate external dependencies such as the mapping API and network responses. This allows testing of edge cases and failure scenarios without relying on external services.
	
	\textbf{Lifecycle testing} will evaluate the system across the full game lifecycle, including game creation, player joining, gameplay progression, and game termination, ensuring stability and correctness throughout.
	
	\textbf{Integration testing} will evaluate how different system components interact, particularly communication between the frontend, backend, and mapping API. These tests will ensure that player actions are correctly transmitted to the server, processed by the game engine, and synchronised across all connected clients. While initially performed manually, automation may be introduced if time permits.
	
	\textbf{Functional testing} will validate that the full gameplay experience follows the rules of Scotland Yard. Test scenarios will simulate different game situations to confirm that movement rules, hidden player mechanics, ticket restrictions, and win conditions are implemented correctly.
	
	\textbf{Performance testing} will evaluate the responsiveness and stability of the multiplayer system. This will involve measuring the latency of game updates, testing synchronisation between multiple clients, and verifying that the server can maintain a consistent game state under concurrent player interactions.
	
	To ensure the quality of the testing process itself, \textbf{meta-testing techniques} will be applied. These include measuring \textbf{test coverage} to ensure sufficient code is exercised, \textbf{mutation testing} to evaluate the effectiveness of test cases, \textbf{fuzz testing} to uncover unexpected edge cases through randomised inputs, and \textbf{property-based testing} to validate general invariants of the game logic.
	
	\textbf{User evaluation} will be conducted to assess usability and gameplay experience. Participants will play the game using the browser interface and provide feedback through the System Usability Scale (SUS) questionnaire, a widely used standardised usability metric. Additional qualitative feedback will be collected to identify interface issues, clarity of game information, and overall user satisfaction. The results will be used to guide improvements to the user interface and interaction design.
	
	Together, these evaluation methods will demonstrate that the system correctly implements the game rules, performs reliably in a multiplayer environment, and provides an effective and usable gameplay experience.
	
	\section{Resourcing and Ethics}	
	\subsection{Hardware}
	
	There aren't many hardware needs for this project. The main thing other than a laptop would be a server to host the project once its completed. I have a home server that I'm planning to host this on permanently once its finished so theres no extra costs.
	
	\subsection{Software, Datasets and Models}
		
	The specialised software and resources planned for this project include:
	
	\begin{itemize}
		\item \textbf{Frontend framework:} Vue.js or React.js for building the user interface
		\item \textbf{Mapping API:} for rendering the Wellington region and game locations
		\item \textbf{Backend language:} Java or JavaScript for implementing game logic and multiplayer communication
		\item \textbf{Deployment and infrastructure:}
		\begin{itemize}
			\item \textbf{Hosting hardware:} A physical device for deployment, using a personal home server
			\item \textbf{Container hosting platform:} Virtualisation or cloud platforms such as Proxmox, AWS, or Azure
			\item \textbf{Containerisation:} Docker for packaging and deploying application components
			\item \textbf{Networking and security:} Cloudflare for HTTPS and secure external access
		\end{itemize}
	\end{itemize}
	
	\subsection{Space, Virtual Resources and Access}
	
	The project will make use of both physical and virtual resources to support development, testing, and deployment:
	
	\begin{itemize}
		\item \textbf{Physical resources:}
		\begin{itemize}
			\item Home workspace for development and testing
			\item Personal computer and network access for local multiplayer testing
		\end{itemize}
		
		\item \textbf{Virtual resources:}
		\begin{itemize}
			\item Virtual machines hosted on Proxmox to isolate backend and frontend services
			\item Docker containers for deploying the application in a controlled environment
			\item Cloudflare services for external network access and testing
			\item Online mapping API services to provide geographic data and map rendering
		\end{itemize}
	\end{itemize}
	
	Access to these resources will be through personal equipment and accounts, with virtualised environments configured on the home server. This setup allows for controlled testing of multiplayer interactions, deployment scenarios, and integration with external services while avoiding the need for specialised labs or external facilities.
	
	\subsection{Budget}
	
	\begin{itemize}
		\item Physical Scotland Yard boardgame (for reference and testing) supervisor already ordered: \$80
		\item Claude Code subscription 6 months approved by supervisor: \$204
		\item Domain name for the website: \$10
		\item Some mapping API, the ideal is to find a free api but there could potentually be a cost here.
	\end{itemize}
	
	Total: \$294
	
	\subsection{Ethics}
	
	No additional ethical considerations are anticipated. Testing will be conducted primarily through software-based methods, and any user testing will be carried out with voluntary participation from students. Ethics for testing on fellow students has been approved by the course coordinator already.
	
	\subsection{Safety}
	
	Standard safety precautions for computer-based work will be observed, including managing eye strain, taking regular breaks to avoid prolonged sitting, and maintaining ergonomic posture to reduce the risk of wrist or back discomfort.
	
	\subsection{Intellectual Property}
	
	Intellectual property resulting from this project will be owned by the author, Mingzhe Bao. The university will retain permission to use the project materials for promotional purposes, including Open Day events and advertising.
	
	Permission will be sought from Ravensburger to use the \textit{Scotland Yard} intellectual property as the basis for this project. If permission is not granted, the project will be adapted to use an original theme and modified game elements that preserve the core mechanics while avoiding the use of protected intellectual property.
	
\end{document}
